\chapter[Lezione 20 --- 18 novembre 2026]{Lezione 20: limiti trigonometrici e logaritmici}
\label{chap:lezione-20}

\section{Disuguaglianze trigonometriche}
\label{sec:lez20-disuguaglianze-trigonometriche}

\begin{proposition}[Disuguaglianze per il seno]
\label{prop:lez20-disuguaglianze-seno}
Per ogni \(x\in(0,\pi/2)\), si ha
\[
  0<\sin x<x<\tan x.
\]
Inoltre, per ogni \(x\in\R\), vale
\[
  |\sin x|\le|x|.
\]
\end{proposition}

\begin{proof}
\emph{Primo punto.}
Si consideri la circonferenza goniometrica e si indichino
con \(O\) il centro, \(B=(1,0)\), \(C=(\cos x,\sin x)\)
e \(D=(1,\tan x)\). Per \(0<x<\pi/2\), il triangolo
\(OCB\) è contenuto nel settore circolare \(OCB\),
a sua volta contenuto nel triangolo \(OBD\).

\begin{figure}[htbp]
  \centering
  \begin{tikzpicture}[unipd/diagram,scale=2.8]
    \def\ang{40}
    \coordinate (O) at (0,0);
    \coordinate (B) at (1,0);
    \coordinate (C) at ({cos(\ang)},{sin(\ang)});
    \coordinate (A) at ({cos(\ang)},0);
    \coordinate (D) at (1,{tan(\ang)});
    \draw[->,unipd-muted] (-0.10,0)--(1.32,0) node[right] {$x$};
    \draw[->,unipd-muted] (0,-0.10)--(0,1.10) node[above] {$y$};
    \fill[unipd-primary,opacity=.10]
      (O)--(B) arc[start angle=0,end angle=\ang,radius=1]--cycle;
    \draw[unipd-primary,very thick]
      (B) arc[start angle=0,end angle=\ang,radius=1];
    \draw[unipd-primary,line width=.8pt] (O)--(C)--(B);
    \draw[unipd-muted,densely dashed] (C)--(A);
    \draw[unipd-muted,line width=.8pt] (B)--(D);
    \draw[unipd-primary,line width=.8pt] (O)--(D);
    \draw[unipd-muted] (0.20,0) arc[start angle=0,end angle=\ang,radius=.20];
    \node[below left] at (O) {$O$};
    \node[below] at (A) {$A$};
    \node[below] at (B) {$B$};
    \node[above left] at (C) {$C$};
    \node[above right] at (D) {$D$};
    \node at (.29,.11) {$x$};
  \end{tikzpicture}
  \caption{Confronto tra il triangolo \(OCB\), il settore
  circolare \(OCB\) e il triangolo \(OBD\) nella
  circonferenza di raggio unitario.}
  \label{fig:lez20-confronto-aree}
\end{figure}

Poiché il raggio è unitario, le rispettive aree sono
\[
\begin{aligned}
  \operatorname{area}(\triangle OCB)&=\frac{\sin x}{2},\\
  \operatorname{area}(\text{settore }OCB)&=\frac{x}{2},\\
  \operatorname{area}(\triangle OBD)&=\frac{\tan x}{2}.
\end{aligned}
\]
Dalla stretta inclusione tra le tre regioni segue
\[
  \frac{\sin x}{2}<\frac{x}{2}<\frac{\tan x}{2}.
\]
Poiché \(\sin x>0\), si conclude
\[
  0<\sin x<x<\tan x.
\]

\emph{Secondo punto.}
Per \(0<x<\pi/2\), dal primo punto segue
\(|\sin x|<|x|\). Se \(x\ge\pi/2\), invece,
\[
  |\sin x|\le1<\frac\pi2\le x.
\]
Per \(x<0\), posto \(u:=-x>0\),
dalla disparità del seno segue
\[
  |\sin x|=|\sin u|\le u=|x|.
\]
Infine per \(x=0\) si ha \(|\sin0|=0\).
Pertanto la disuguaglianza vale per ogni \(x\in\R\).
\end{proof}

\section{Limiti delle funzioni trigonometriche}
\label{sec:lez20-limiti-trigonometrici}

\begin{proposition}[Limiti di seno, coseno, tangente e cotangente]
\label{prop:lez20-limiti-trigonometrici}
Per ogni \(x_0\in\R\), si ha
\[
  \lim_{x\to x_0}\sin x=\sin x_0,
  \qquad
  \lim_{x\to x_0}\cos x=\cos x_0.
\]
Inoltre,
\begin{enumerate}
  \item se \(x_0\ne\pi/2+k\pi\) per ogni \(k\in\Z\),
  \[
    \lim_{x\to x_0}\tan x=\tan x_0;
  \]
  \item se \(x_0\ne k\pi\) per ogni \(k\in\Z\),
  \[
    \lim_{x\to x_0}\cot x=\cot x_0.
  \]
\end{enumerate}
\end{proposition}

\begin{proof}
\emph{Primo punto: seno.}
Per le formule di prostaferesi e la
\cref{prop:lez20-disuguaglianze-seno}, si ha
\[
  |\sin x-\sin x_0|
  =2\left|\cos\frac{x+x_0}{2}\right|
    \left|\sin\frac{x-x_0}{2}\right|
  \le2\left|\frac{x-x_0}{2}\right|
  =|x-x_0|.
\]
Poiché \(|x-x_0|\to0\) per \(x\to x_0\), dal
\cref{thm:lez12-tre-funzioni} segue
\[
  \lim_{x\to x_0}|\sin x-\sin x_0|=0,
\]
e quindi \(\lim_{x\to x_0}\sin x=\sin x_0\).

\emph{Secondo punto: coseno.}
Per le formule di prostaferesi e la
\cref{prop:lez20-disuguaglianze-seno}, si ha
\[
  |\cos x-\cos x_0|
  =2\left|\sin\frac{x+x_0}{2}\right|
    \left|\sin\frac{x-x_0}{2}\right|
  \le2\left|\frac{x-x_0}{2}\right|
  =|x-x_0|.
\]
Poiché \(|x-x_0|\to0\) per \(x\to x_0\), dal
\cref{thm:lez12-tre-funzioni} segue
\[
  \lim_{x\to x_0}|\cos x-\cos x_0|=0,
\]
e quindi \(\lim_{x\to x_0}\cos x=\cos x_0\).

\emph{Terzo e quarto punto.}
Per ogni \(x\) nei rispettivi domini,
\[
  \tan x=\frac{\sin x}{\cos x},
  \qquad
  \cot x=\frac{\cos x}{\sin x}.
\]
Se \(\cos x_0\ne0\), rispettivamente
\(\sin x_0\ne0\), i limiti appena dimostrati
e il \cref{thm:lez14-limite-rapporto} danno
\[
  \lim_{x\to x_0}\tan x
    =\frac{\sin x_0}{\cos x_0}=\tan x_0,
  \qquad
  \lim_{x\to x_0}\cot x
    =\frac{\cos x_0}{\sin x_0}=\cot x_0,
\]
con ciascuna uguaglianza intesa nel proprio dominio.
\end{proof}

\clearpage

\section{Limiti notevoli trigonometrici}
\label{sec:lez20-limiti-notevoli-trigonometrici}

Si dimostrano il settimo e l'ottavo punto del
\cref{thm:lez17-limiti-notevoli}, utilizzando
le disuguaglianze e i limiti precedenti.

\begin{theorem}[Limite notevole del seno]
\label{thm:lez20-limite-seno}
Si ha
\[
  \lim_{x\to0}\frac{\sin x}{x}=1.
\]
\end{theorem}

\begin{proof}
Si consideri dapprima \(x\to0^+\).
Per \(x\in(0,\pi/2)\), dalla
\cref{prop:lez20-disuguaglianze-seno} segue
\[
  0<\sin x<x<\frac{\sin x}{\cos x}.
\]
Dividendo per \(\sin x>0\), si ottiene
\[
  1<\frac{x}{\sin x}<\frac1{\cos x}.
\]
Per la \cref{prop:lez20-limiti-trigonometrici},
\[
  \lim_{x\to0^+}\cos x=1.
\]
Dal \cref{thm:lez12-tre-funzioni} si ricava
\[
  \lim_{x\to0^+}\frac{x}{\sin x}=1,
\]
e dal \cref{thm:lez14-limite-rapporto} segue
\[
  \lim_{x\to0^+}\frac{\sin x}{x}=1.
\]
La funzione \(x\mapsto\sin x/x\) è pari nel
suo dominio, poiché sia il seno sia la funzione
\(x\mapsto x\) sono dispari:
\[
  \frac{\sin(-x)}{-x}
  =\frac{-\sin x}{-x}
  =\frac{\sin x}{x}.
\]
Per la \cref{prop:lez19-limiti-funzione-pari},
anche il limite sinistro vale \(1\). I due limiti
unilateri coincidono, quindi per il
\cref{thm:lez14-limite-unilateri} si conclude
\[
  \lim_{x\to0}\frac{\sin x}{x}=1.
\]
\end{proof}

Il limite afferma che, in prossimità di zero,
\(\sin x\) e \(x\) sono asintoticamente equivalenti:
il loro rapporto tende a \(1\). Pertanto, nel
calcolo dei limiti, possono essere sostituiti
l'uno con l'altro nei prodotti e nei quozienti.

\begin{proposition}[Limite notevole del coseno]
\label{prop:lez20-limite-coseno}
Si ha
\[
  \lim_{x\to0}\frac{1-\cos x}{x^2}=\frac12.
\]
\end{proposition}

\begin{proof}
Per \(x\in(-\pi,\pi)\setminus\{0\}\),
si moltiplicano numeratore e denominatore
per \(1+\cos x\):
\[
\begin{aligned}
  \frac{1-\cos x}{x^2}
  &=\frac{1-\cos^2 x}{x^2(1+\cos x)}
  &=\left(\frac{\sin x}{x}\right)^2
    \frac1{1+\cos x}.
\end{aligned}
\]
Per il \cref{thm:lez20-limite-seno}
e la \cref{prop:lez20-limiti-trigonometrici},
applicando i \cref{thm:lez14-limite-prodotto,thm:lez14-limite-rapporto},
si ottiene
\[
  \lim_{x\to0}\frac{1-\cos x}{x^2}
  =1^2\cdot\frac1{1+1}=\frac12.
\]
\end{proof}

\clearpage

\section{Limiti delle funzioni logaritmiche}
\label{sec:lez20-limiti-logaritmo}

\begin{proposition}[Limiti del logaritmo]
\label{prop:lez20-limiti-logaritmo}
Siano \(a>0\), con \(a\ne1\), e \(x_0>0\). Allora
\[
  \lim_{x\to1}\log_a x=0,
  \qquad
  \lim_{x\to x_0}\log_a x=\log_a x_0.
\]
\end{proposition}

\begin{proof}
Si dimostra anzitutto il limite in \(x_0=1\).
Sia \(\varepsilon>0\).

\emph{Primo caso: \(a>1\).}
Poiché il logaritmo di base \(a\) è strettamente
crescente, per la \cref{ex:lez10-logaritmo-monotono}
si ha
\[
\begin{aligned}
  |\log_a x|<\varepsilon
  &\quad\Longleftrightarrow\quad
    -\varepsilon<\log_a x<\varepsilon
  &\quad\Longleftrightarrow\quad
    a^{-\varepsilon}<x<a^\varepsilon.
\end{aligned}
\]
Poiché \(a^{-\varepsilon}<1<a^\varepsilon\),
si scelga
\[
  \delta_\varepsilon
  :=\min\{1-a^{-\varepsilon},a^\varepsilon-1\}>0.
\]
Per ogni \(x>0\) con \(|x-1|<\delta_\varepsilon\),
risulta \(a^{-\varepsilon}<x<a^\varepsilon\),
e quindi \(|\log_a x|<\varepsilon\).

\emph{Secondo caso: \(0<a<1\).}
Ponendo \(b:=1/a>1\), per la formula di cambiamento
di base si ha \(\log_a x=-\log_b x\).
Il logaritmo di base \(a\) è quindi strettamente
decrescente. Pertanto
\[
\begin{aligned}
  |\log_a x|<\varepsilon
  &\quad\Longleftrightarrow\quad
    -\varepsilon<\log_a x<\varepsilon
  &\quad\Longleftrightarrow\quad
    a^\varepsilon<x<a^{-\varepsilon}.
\end{aligned}
\]
Poiché \(a^\varepsilon<1<a^{-\varepsilon}\),
si scelga
\[
  \delta_\varepsilon
  :=\min\{1-a^\varepsilon,a^{-\varepsilon}-1\}>0.
\]
Per ogni \(x>0\) con \(|x-1|<\delta_\varepsilon\),
si ha anche in questo caso
\(|\log_a x|<\varepsilon\).
Per la \cref{def:lez08-limite}, in entrambi i casi
segue
\[
  \lim_{x\to1}\log_a x=0.
\]

Sia ora \(x_0>0\) arbitrario. Ponendo
\(t:=x/x_0\), per \(x\to x_0\) si ha \(t\to1\);
inoltre, se \(x\ne x_0\), allora \(t\ne1\).
Dalle proprietà dei logaritmi segue
\[
  \log_a x=\log_a\left(\frac{x}{x_0}\right)
          +\log_a x_0.
\]
Per il \cref{thm:lez16-limite-sostituzione},
applicato al limite appena dimostrato, e per il
\cref{thm:lez14-limite-somma}, si conclude
\[
  \lim_{x\to x_0}\log_a x
  =0+\log_a x_0=\log_a x_0.
\]
\end{proof}

\begin{exercise}[Un limite con funzioni trigonometriche composte]
\label{ex:lez20-coseno-seno-rapporto}
Studiare il limite
\[
  \lim_{x\to0}\frac{1-\cos(\sin x)}{x^3}.
\]
\end{exercise}

\begin{proof}[Soluzione dell'esercizio]
Per \(x\ne0\) sufficientemente vicino a zero,
si può scrivere
\[
  \frac{1-\cos(\sin x)}{x^3}
  =\frac{1-\cos(\sin x)}{(\sin x)^2}
   \left(\frac{\sin x}{x}\right)^2\frac1x.
\]
Ponendo \(t:=\sin x\), per \(x\to0\)
si ha \(t\to0\), per la
\cref{prop:lez20-limiti-trigonometrici}. Inoltre,
per \(x\ne0\) vicino a zero si ha \(t\ne0\).
Per il \cref{thm:lez16-limite-sostituzione}
e la \cref{prop:lez20-limite-coseno},
\[
  \lim_{x\to0}
  \frac{1-\cos(\sin x)}{(\sin x)^2}=\frac12.
\]
Dal \cref{thm:lez20-limite-seno} segue anche
\[
  \lim_{x\to0}\left(\frac{\sin x}{x}\right)^2=1.
\]
Il prodotto dei primi due fattori tende quindi a
\(1/2>0\). Poiché \(1/x\to+\infty\) per
\(x\to0^+\) e \(1/x\to-\infty\) per \(x\to0^-\),
per il \cref{thm:lez14-limite-prodotto} si ottiene
\[
  \lim_{x\to0^+}\frac{1-\cos(\sin x)}{x^3}=+\infty,
  \qquad
  \lim_{x\to0^-}\frac{1-\cos(\sin x)}{x^3}=-\infty.
\]
I due limiti unilateri sono distinti; pertanto,
per il \cref{thm:lez14-limite-unilateri},
il limite richiesto non esiste.
\end{proof}